\documentclass[11pt,a4paper]{article}
\usepackage[utf8]{inputenc}
\usepackage[T1]{fontenc}
\usepackage{geometry}
\geometry{top=1in, bottom=1in, left=1in, right=1in}
\usepackage{hyperref}
\usepackage{booktabs}
\usepackage{amsmath}
\usepackage{listings}
\usepackage{xcolor}
\usepackage{graphicx}
\usepackage{array}
\usepackage{fancyhdr}
\usepackage{titlesec}
\usepackage{longtable}

\hypersetup{
    colorlinks=true,
    linkcolor=purple,
    filecolor=magenta,      
    urlcolor=cyan,
    pdftitle=Project Genesis Engineering Principles,
    pdfpagemode=FullScreen,
}

\definecolor{codegray}{rgb}{0.96,0.96,0.96}
\definecolor{codeborder}{rgb}{0.85,0.85,0.85}
\definecolor{commentgreen}{rgb}{0.12,0.55,0.22}
\definecolor{keywordblue}{rgb}{0.1,0.2,0.8}

\lstset{
    backgroundcolor=\color{codegray},
    basicstyle=\ttfamily\small,
    breakatwhitespace=false,
    breaklines=true,
    captionpos=b,
    commentstyle=\color{commentgreen},
    frame=single,
    rulecolor=\color{codeborder},
    keepspaces=true,
    keywordstyle=\color{keywordblue},
    language=Python,
    showspaces=false,
    showstringspaces=false,
    showtabs=false,
    tabsize=4
}

\pagestyle{fancy}
\fancyhf{}
\rhead{\small Project Genesis Engineering Principles}
\lhead{\small Project Genesis}
\cfoot{\thepage}

\title{\textbf{Project Genesis Engineering Principles}}
\author{Project Genesis Research Repository}
\date{\small Generated: July 2026}

\begin{document}

\maketitle
\thispagestyle{empty}

\newpage
\section*{Project Genesis Engineering Principles}

This document outlines the core architectural and software engineering principles governing all subsystems of Genesis (the Simulation Engine and the Web Portal).

\noindent\rule{\textwidth}{0.4pt}

\subsection*{1. Single Subsystem Responsibility}
Subsystems are strictly isolated:
\begin{itemize}
  \item \textbackslash{}textbf{Simulation Engine (\textbackslash{}texttt{world/})}: Handles cellular and agent cognitive operations. It does not contain server or database integration logic.
  \item \textbackslash{}textbf{Web Portal (\textbackslash{}texttt{portal/})}: Handles ingestion validation, indexing, and visualization. It does not compute physics simulation steps.
\end{itemize}

\subsection*{2. The Experiment Package is the Contract}
The exported ZIP archive is the absolute source of truth. All subsystems rely on the schema contracts defined in \textbackslash{}texttt{shared/experiment\textbackslash{}\_schema.json}. We never compile database fields or client features that cannot be derived from this package contract.

\subsection*{3. Narrative Prioritized Over Analytics}
The human chronicle of events (Founding, Expansion, Conflict, Decline, Collapse) is the entry gate to every experiment record. We present historical story records first to build emotional and context understanding, keeping data graphs decoupled inside the Observatory.

\subsection*{4. Scientific Reproducibility Above All}
Simulations must be verifiable:
\begin{itemize}
  \item All assets (configs, events, raw populations) are downloadable.
  \item Every experiment record includes citation support (BibTeX, APA) and a provenance tracker showing engine versions and staging logs.
\end{itemize}

\subsection*{5. Unified Schema Control}
To prevent database, backend, and engine drifts, schema specifications are never duplicated. Shared types and JSON schemas live in the root \textbackslash{}texttt{shared/} directory and are referenced directly during build steps.

\subsection*{6. Every Feature Must Answer a Research Goal}
We do not add dashboards, telemetry charts, or spatial tracking features simply because the data exists. Every visualization must act as a target "Laboratory Instrument" resolving a specific research hypothesis.

\subsection*{7. Progressive Enhancement}
We prioritize deep, robust implementations over feature accumulation. We build foundation layers (ingestion validation, transactional rollback, dynamic timeline chronicle) before visual complexity (interactive replays, animations, search graphs).

\end{document}
